%-------------------------------------------------------------
\subsection*{R2V generation model}\label{sec:baseline}
%-------------------------------------------------------------

\textbf{Backbone.}
We develop a multi-control video generation model for the R2V
task, built on
Cosmos-Transfer2.5~\cite{cosmos2025transfer,agarwal2025cosmos},
a rectified-flow diffusion transformer (DiT) pre-trained for
text- and control-conditioned driving video synthesis. In
particular, we adopt the 2B \emph{Auto Multiview} variant
(\texttt{nvidia/Cosmos-Transfer2.5-2B}), which generates
29-frame $1280\!\times\!720$ clips at 10\,Hz across seven
synchronised vehicle cameras. The backbone operates in the
latent space of the Wan~2.1 video
tokenizer~\cite{wan2pt1}, which compresses a
$29\!\times\!720\!\times\!1280$ pixel clip to a latent of
temporal length $T_{\mathrm{lat}}\!=\!8$ with a $4\!\times$
temporal and $16\!\times\!16$ spatial compression ratio, and
reconstructs pixel frames through its paired decoder. Text is
encoded by Cosmos Reason1.1-7B~\cite{cosmos_reason} using the
full-concat embedding strategy. Cross-view consistency over
the seven surround-view cameras is provided by the multi-view
variant of the DiT, which interleaves spatial, temporal and
cross-view self-attention within every block and augments
token positions with a camera-index
embedding~\cite{cosmos2025transfer}.

\textbf{Control-unit injection.}
The three channels of the R2V control unit
(Methods: \emph{R2V control unit}) are injected into the
backbone through a VACE-style sidecar~\cite{jiang2025vace}.
Each of the three condition videos is independently encoded
by the frozen Wan~2.1 VAE to a 16-channel latent, and the
three latents are concatenated along the channel dimension
into a single 48-channel control latent,
%
\begin{equation}
  \label{eq:ctrl_cat}
  \mathbf{c}_{\mathrm{R2V}} =
  \big[\,\mathbf{z}_{\mathrm{layout}}\,;\,\mathbf{z}_{\mathrm{photo}}\,;\,\mathbf{z}_{\mathrm{geo}}\,\big]
  \in \mathbb{R}^{48\times V T_{\mathrm{lat}}\times H'\times W'},
\end{equation}
%
where $V\!=\!7$ denotes the number of surround-view cameras,
$T_{\mathrm{lat}}\!=\!8$ the latent temporal length, and
$H'\!=\!45$, $W'\!=\!80$ the spatial dimensions of the latent
tensor for a $1280\!\times\!720$ pixel input; the layout,
photometric and geometric latents occupy channels 0--15,
16--31 and 32--47, respectively. A \emph{control embedder}
patchifies $\mathbf{c}_{\mathrm{R2V}}$ into tokens whose
length matches that of the DiT token stream, and an optional
\emph{input hint block} refines the embedded tokens with a
per-frame multilayer perceptron. The tokens are then
processed by a stack of lightweight \emph{control
blocks}---each sharing the attention and feed-forward
structure of a base DiT block---inserted at every seventh
layer of the backbone (layer indices $\{0,7,14,\ldots\}$).
Every control block contributes an additive residual
modulation to its companion base block through
zero-initialised linear projections, so that at
initialisation the backbone output is identical to the
unconditioned one and the control pathway is learnt purely
from subsequent gradient updates. Aside from the
channel-count adjustment of the control embedder, no
architectural modification is introduced relative to the
public Cosmos-Transfer2.5 release.

\textbf{Channel assignment and weight reuse.}
The channel ordering of the R2V control unit is deliberate.
The publicly released Cosmos-Transfer2.5-2B Auto Multiview
checkpoint was trained with a single control modality
(\texttt{hdmap\_bbox})~\cite{cosmos2025transfer} that
occupies the first 16 channels of its control embedder. By
assigning the layout channel---whose pixel statistics (solid
polygons of categorical colours) most closely match the
\texttt{hdmap\_bbox} prior---to those 16 channels, we inherit
the pre-trained weights and amortise the cost of learning a
semantic-layout control pathway from scratch. The remaining
32 channels, which encode the photometric and geometric
signals, are initialised randomly and specialised on our
dataset during post-training. The unified interface avoids
instantiating three parallel ControlNet branches and, by
construction, allows the three complementary signals to be
attended to jointly within every control block.

\textbf{Text conditioning.}
Because the dataset is collected at a fixed intersection,
clip-level narration carries little information beyond what
the R2V control unit already encodes, while dense captions
risk overfitting the model to templated surface forms. We
therefore adopt a structured scene-level prompt composed of
(i) a view-identifier phrase (e.g.\ ``Front wide view.''),
(ii) an ego-motion direction phrase derived from the recorded
trajectory (e.g.\ ``The ego vehicle is travelling from west
to east.''), and (iii) a fixed intersection description
covering the geographic and seasonal context, road
infrastructure and typical traffic participants. This
decomposition delegates dynamic information to the control
unit and static scene context to the text prompt. During
training the prompt is dropped with probability 0.2 to enable
classifier-free guidance at inference time.

\textbf{Trainable parameters.}
Following the context-adapter tuning strategy of
VACE~\cite{jiang2025vace}, the entire pre-trained backbone is
kept frozen and only the control embedder, the control blocks
and the input hint block receive gradients during
post-training. Trainable parameters account for
$3.29\!\times\!10^{8}$ ($13.8\%$) of the
$2.39\!\times\!10^{9}$ parameters of the base network,
concentrated in the control blocks ($2.98\!\times\!10^{8}$;
$90.6\%$ of trainable parameters), with the input hint block
($2.94\!\times\!10^{7}$; $8.9\%$) and the control embedder
($1.6\!\times\!10^{6}$; $0.5\%$) contributing the remainder.
